
\section*{Executive Summary}
In recent decades, the surge of digital technologies has reshaped the way data are collected, stored, and used. This chapter charts a historical and conceptual journey through the transformation from raw data to actionable knowledge, highlighting key shifts in information systems. Section~1.1 introduces the theme of \emph{digitalization}, showing how organizations moved from manual record-keeping to modern data-driven practices. It reviews foundational concepts such as the DIKW (Data, Information, Knowledge, Wisdom) pyramid and the emergence of \emph{business intelligence} (BI) and knowledge discovery methodologies (e.g., CRISP-DM) that formalized the process of turning data into insight. Section~1.2 then broadens the perspective to the societal level, examining how \emph{IoT} and \emph{Industry 4.0} (smart manufacturing) have multiplied data sources, necessitating new storage and processing paradigms. We discuss the limitations of traditional relational databases and data warehouses when confronted with \emph{Big Data} volumes, and the rise of NoSQL, data lakes, and lakehouse architectures to address these challenges. Section~1.3 tackles the current and future data economy: it explores how data have become a critical resource (“the new oil” metaphor), and how emerging technologies such as \emph{cloud platforms}, \emph{artificial intelligence (AI)} and \emph{foundation models} (large-scale language/image models) further elevate the importance of data. We also examine concepts like \emph{Digital Twins} and autonomous AI agents as exemplars of the intertwining of data and decision-making. Throughout, we emphasize causal motivations for each evolution (e.g. scalability, real-time needs) and conclude by identifying open research challenges arising from this data-centric world. This cohesive narrative sets the stage for the rest of the thesis by illustrating why data architectures and methods (from lakes to twins) are both crucial and timely research topics.

\clearpage

\section{Digitalization: From Data to Knowledge}
Digitalization refers to the ongoing transformation of processes, organizations and societies through information technology. In its broadest sense, it signifies a shift from analog, manual operations to the systematic use of digital data and computational models.  Historically, businesses accumulated more data than they could effectively exploit. Early efforts in the 1990s led to the development of \emph{data warehouses} and \emph{business intelligence} (BI) systems, which aimed to consolidate transactional data into a structured repository for reporting and decision support.  Influential works by Davenport and others (e.g. “\emph{Competing on Analytics}”, 2006) argued that competitive advantage increasingly depended not on products alone but on the ability to analyze vast data stores and act on insights. This perspective helped promote the cultural shift toward valuing data-driven decision-making throughout organizations.

A key conceptual framework for understanding this shift is the DIKW hierarchy (Data--Information--Knowledge--Wisdom). In this model (attributed to Ackoff, 1989), \emph{data} are raw facts or observations, \emph{information} is data that have been organized or given context, \emph{knowledge} is information that has been processed or interpreted, and \emph{wisdom} is the prudent application of knowledge. The DIKW pyramid illustrates how incremental processing of data yields higher-value insights (see Figure). Researchers and practitioners have used this hierarchy to justify investments in transforming data into actionable knowledge. For example, data by itself has limited utility; it must be cleaned, integrated, and analyzed to become meaningful information (e.g. a summary report or visualization), which then can inform decisions.

\subsection{The Digital Transformation}
The process of digital transformation is inherently evolutionary. Initially, many organizations digitized routine record-keeping (e.g. inventory, transactions) using relational databases. Over time, as more aspects of operations became automated, the volume and variety of data grew. For instance, sensor networks and machine logs began generating high-frequency data streams. To make sense of this data deluge, companies established data warehouses where data from diverse operational systems were extracted, transformed, and loaded (ETL) into a centralized schema. This enabled reporting and batch analysis, a hallmark of classic BI systems. Crucially, top management often championed BI initiatives: Davenport (2006) noted that firms “competing on analytics” treat data as a strategic asset, mandating organization-wide analytics programs. 

As data quantities outpaced the ability of traditional systems, more sophisticated \emph{data-driven} methodologies emerged. Rather than just examining past trends, organizations wanted predictive and prescriptive insights. This led to \emph{knowledge discovery} and data mining, formalized in works like Fayyad et al. (1996). The Knowledge Discovery in Databases (KDD) process defined steps from raw data to useful knowledge, including selection, preprocessing, transformation, mining, and interpretation. Likewise, the Cross-Industry Standard Process for Data Mining (CRISP-DM) introduced around 2000 provided a six-stage cycle (business understanding, data understanding, data preparation, modeling, evaluation, deployment) to guide analytics projects. Together, these frameworks underscored that converting data into knowledge is not a one-time task but a cyclical process involving both technical and business understanding.

\subsection{The DIKW Hierarchy and Business Intelligence}
Under the DIKW framework, earlier stages of data management correspond to moving from the bottom to the middle of the pyramid. In practice, \emph{information} is produced by summarizing or visualizing data (e.g. charts of sales over time), often through BI dashboards. These dashboards aggregate operational data (OLTP) into summaries that managers can interpret. But information alone is not enough; organizations realized the need to generate \emph{knowledge} – meaning patterns or rules extracted from information. Early knowledge discovery efforts applied statistical analysis and machine learning to identify correlations and trends (for example, customer segmentation or predictive risk models).

Business Intelligence played a central role in this stage of digitalization. BI platforms integrate data from multiple sources and enable interactive exploration. They often rely on data warehouses where data is cleansed and modeled (using star schemas, as advocated by Kimball or Inmon) to support fast query performance. The evolution from manual reports to BI can be seen as the transformation of \emph{data into knowledge}: once data is cleaned and modeled, users can pose queries that yield informative patterns. Studies have shown that BI adoption correlates with improved decision-making capability and competitive performance (e.g. Davenport 2006).

\subsection{Knowledge Discovery and CRISP-DM}
Building on BI, the late 1990s and early 2000s saw the formalization of \emph{data mining} and knowledge discovery processes. Fayyad et al. (1996) defined KDD as the non-trivial process of identifying valid, novel, useful, and understandable patterns in data. The data mining step within KDD uses algorithms to find predictive or descriptive patterns (classification, clustering, association rules, etc.). CRISP-DM (Chapman et al., 2000) emerged as a pragmatic standard: it emphasized understanding the business context and iterative refinement. In CRISP-DM’s cycle, one starts with a business problem, examines the data, prepares it (cleaning, feature engineering), applies models, and then evaluates and deploys the results. These methodologies underscored that achieving “knowledge” (and eventually “wisdom”) from data requires disciplined procedures, including defining objectives, iteratively modeling, and embedding the output back into organizational decisions.

CRISP-DM and KDD also foreshadowed the ongoing cycle of digital transformation: as new data sources emerge (e.g. web logs, sensors), the data preparation and modeling tasks restart. For example, discovering that sensor readings correlate with machine failures leads to a predictive maintenance model, which in turn becomes another form of \emph{wisdom} when used to proactively schedule repairs. Through these processes, data evolves from inert records into actionable knowledge that can optimize operations, reduce costs, and enable new business models.

\section{The Role of Data in a Digital Society}
As society digitalized, the pervasiveness of data extended beyond individual enterprises to entire systems and networks. Two interlinked trends are especially noteworthy: the proliferation of connected devices (the Internet of Things, IoT) and the vision of \emph{Industry 4.0} (smart factories, cyber-physical systems). These developments turned previously isolated data silos into streams of real-time information. Consequently, the IT infrastructure evolved from standalone databases to distributed, scalable platforms capable of handling \emph{Big Data}.

\subsection{The Internet of Things (IoT) and Cyber-Physical Systems}
The Internet of Things (IoT) refers to networks of physical objects (“things”) embedded with sensors, actuators, and communication interfaces. Examples range from household appliances and wearable devices to industrial sensors on a factory floor. In essence, IoT connects the physical world to the digital: machines and sensors continuously generate data about environmental and operational conditions (temperature, vibration, usage, etc.). This sensor data is then transmitted (often via wireless networks) to centralized or cloud-based processing systems. 

In industry, IoT combines with cyber-physical systems (CPS) to form smart manufacturing environments. IoT devices on equipment capture data, while CPS provides computation and control tightly integrated with physical processes. For instance, a wind turbine may have IoT sensors reporting blade stress and a CPS adjusting blade pitch in real-time. The integration of IoT means that traditional manufacturing (Industry 3.0) has become the cyber-physical factories of Industry 4.0: highly instrumented, networked, and automated.

A key consequence of IoT is an explosion of data. Each sensor can produce readings at high frequency (often multiple times per second). Aggregated across thousands of devices, this quickly becomes a \emph{data stream} that far exceeds what traditional databases were designed for. Moreover, IoT data is often unstructured or semi-structured, and may arrive asynchronously. These characteristics – high volume, velocity, and variety – prompted the need for new data platforms.

\subsection{Industry 4.0 and Smart Manufacturing}
The term \emph{Industry 4.0} (originating around 2013) captures the vision of the fourth industrial revolution: leveraging digital technologies to create \emph{smart factories}. In smart factories, IoT sensors, robotics, cloud computing, and AI work together to enable real-time visibility, self-optimizing processes, and flexible production. As IBM explains, Industry 4.0 is “synonymous with smart manufacturing” and involves the digital transformation of production to achieve real-time decision making, higher productivity, and new levels of customization. 

In practice, Industry 4.0 materializes when data collected from shop-floor devices is combined with enterprise data (e.g. ERP, supply chain, CRM) to give comprehensive insight. For example, data from machines can be linked with maintenance logs to predict breakdowns, or with production schedules to optimize throughput. Crucially, these use cases require data analytics and machine learning on unprecedented scales. Thus, the data-centric nature of Industry 4.0 motivates the second part of our chapter: how data infrastructure evolved to meet this challenge.

\subsection{The Big Data Revolution}
The confluence of IoT, web logs, mobile devices, and social data led to the \emph{Big Data} phenomenon. By the 2010s, organizations were dealing with petabytes of data arriving rapidly. Traditional relational database management systems (RDBMS) became insufficient due to several reasons: scalability limits (vertical scaling only), rigid schemas (difficulty handling unstructured data), and inefficiency for analytic workloads. In response, new data paradigms emerged:

- **NoSQL databases:** These include document stores (e.g. MongoDB), key-value stores (e.g. Redis), wide-column stores (e.g. Cassandra), and graph databases. NoSQL databases relax the rigid schema requirement of SQL, allowing heterogeneous data to coexist. They also typically provide horizontal scalability across distributed clusters, trading some consistency (often offering eventual consistency) for higher throughput. The adoption of NoSQL began in the early 2010s as companies like Google and Facebook needed to manage web-scale workloads.

- **Distributed File Systems and Data Lakes:** Instead of loading data into a structured warehouse, some organizations opted to dump raw data into a storage cluster (often Hadoop HDFS or cloud object storage). This gave rise to the concept of a \emph{data lake}: a unified repository where all data (structured, semi-structured, and unstructured) can be stored at any scale. The data lake allows schema-on-read: data is ingested in its raw form, and schemas or transformations are applied when needed for analysis. While very flexible, data lakes often led to issues of data governance and “data swamps” if not managed properly.

- **Data Warehouses and Federated Architectures:** Parallel to these trends, data warehouses did not disappear; instead, they evolved. Traditional on-premises warehouses (e.g. legacy Teradata or Oracle systems) have largely given way to cloud-based warehouses (e.g. Snowflake, Amazon Redshift). Furthermore, concepts like \emph{data virtualization} and \emph{federated query engines} enable integrating multiple data stores.

Each of these shifts was driven by the need to handle bigger, more complex data. In particular, businesses recognized that relevant data was no longer only in neatly structured tables; it was also in logs, documents, images, and streaming sensor readings. This reality demanded an ecosystem of tools.

\subsection{From Data Warehouses to Data Lakes and Lakehouses}
Data warehouses and data lakes represent two ends of the storage spectrum. A traditional data warehouse is built for structured data and fast analytical queries. It uses ETL pipelines to clean and transform data into a unified schema before loading (schema-on-write). This ensures consistent data quality but requires up-front effort. Data lakes, conversely, prioritize agility and scale: data is ingested first and schema is applied later (schema-on-read). This allows diverse data to be stored cheaply and in large volumes, but can make querying and governance more complex.

Recently, the concept of a \emph{data lakehouse} has emerged as a hybrid architecture. A lakehouse attempts to combine the best of both worlds: it supports ACID transactions and structured tables (like a warehouse) on top of cheap storage (like a lake). For instance, Delta Lake and Apache Iceberg provide metadata layers that enable fast queries and versioning on top of object storage. As Harby and Zulkernine (2025) note, "the data lakehouse merges the strengths of data warehouses and data lakes," allowing quick processing of analytics and the ingestion of high-speed unstructured data. Table summarizes key differences:

This evolution has enabled organizations to ask new kinds of questions. For example, log data from IoT devices can be ingested into a lakehouse, where it can be joined with sales or maintenance data in analytical queries. The lakehouse thus supports both big data machine learning workflows and BI reporting in one platform.

\subsection{Enterprise Data Platforms and Analytics}
Beyond storage, the orchestration of data and analytics has become an industry in itself. Companies now deploy comprehensive \emph{data platforms} (such as AWS, Azure, or Google Cloud offerings) that integrate data ingestion, processing engines, analytics, and governance tools. These platforms often include streaming services (Kafka, Pulsar), batch engines (Spark, Flink), and managed databases. 

The rise of these platforms reflects the recognition that data fuels all decision-support systems (DSS) and advanced applications. Data mining, once a specialized activity, is now embedded as a service: automated anomaly detection, real-time dashboards, and even \emph{Customer Success Management} (CSM) systems rely on feeding large-scale data into AI models. Decision support systems in enterprise often present a single pane of glass combining KPIs, alerts, and predictive insights. In essence, the “role of data” today is foundational: almost every business function (marketing, finance, operations) builds upon data feeds for optimization. This centrality of data has substantial economic and social implications, as we discuss next.

\section{Data is the New Oil}
The phrase “data is the new oil” has become a mantra in popular and business discourse. First uttered by the World Economic Forum in 2011, it emphasizes that data, like oil in the past century, is a valuable raw resource for the modern economy. Indeed, raw data by itself may have limited worth until processed, much like crude oil must be refined before use. However, as many analysts point out, the analogy has limits: data is not depleted by use, and its processing needs are far more varied. Nonetheless, the metaphor captures the explosive growth in data generation and the sense that data fuels innovation.

\subsection{The Data Economy and Strategic Value}
In economic terms, data has moved from a byproduct to a central asset. Companies such as Google, Facebook, Amazon and China’s Tencent have built fortunes on extracting insights from user data. In industry, data-driven firms are often valued highly because of their proprietary datasets and the intelligence gleaned from them. The World Economic Forum noted that data was becoming one of the most important drivers of the Fourth Industrial Revolution. This has led to entire new sectors: data marketplaces, analytics consultancies, and AI service providers.

From a technical perspective, the “new oil” paradigm implies that businesses must invest in the “refinery” and distribution network for data. This means data governance (privacy, security), data pipelines, data scientists, and analytics tools. It also means data centers: companies are building hyperscale data centers at a record pace. For example, investment in data center capacity is booming as AI workloads require specialized hardware (GPUs, TPUs) and vast energy. Projections (IEA, CarbonBrief) suggest data center electricity demand could double by 2030, driven largely by AI. Although data centers currently contribute a small percentage of global emissions (around 1\% by 2024), their growth raises concerns about energy sustainability.

\subsection{Artificial Intelligence and Foundation Models}
The era of Big Data has been paralleled by an AI revolution. In particular, the development of \emph{foundation models} (large-scale AI models pre-trained on massive data) is noteworthy. Foundation models like GPT and BERT are trained on diverse data (text, images, code) to capture broad patterns, and they can be fine-tuned for myriad tasks. Bommasani et al. (2021) define a foundation model as one that is trained on broad data at scale and adaptable to a wide range of downstream tasks. The practical upshot is that these models magnify the value of data: more and better data enables more powerful and generalizable AI.

In industry, foundation models have moved AI from toy demos to mission-critical applications. For example, a language model like GPT can generate business reports, a vision model can inspect products on an assembly line, and multi-modal models can integrate sensor data with natural language. These capabilities rely on the underlying data platform: vast corpora to train on, and real-time data feeds to adapt to new conditions. Hence, the strategic importance of data continues to grow with AI.

\subsection{Digital Twins and Autonomous Agents}
Two emerging paradigms epitomize the fusion of data with operational intelligence: digital twins and autonomous AI agents. A \emph{digital twin} is a virtual replica of a physical system (machine, factory, even city) that is continuously synchronized via data. Pioneered in manufacturing contexts, digital twins enable simulation, optimization, and real-time monitoring of physical assets. For instance, a jet engine’s digital twin uses sensor data to predict wear and schedule maintenance. The rise of digital twins in smart manufacturing and healthcare illustrates how high-fidelity data models become tools for decision-making (thus converting data to knowledge and wisdom).

Autonomous AI \emph{agents} are another frontier. Whereas traditional software executes pre-defined rules, intelligent agents observe an environment and take actions to achieve goals, often with limited human intervention. In business contexts, agents may automate complex tasks (e.g. negotiating trades, coordinating logistics) by using underlying data streams to plan and adapt. Modern agentic systems often leverage foundation models as their reasoning core, integrating external tools and memories. The development of agents reflects a trend towards automating the full loop: from sensing (data acquisition) through reasoning (AI models) to acting (execution), forming a closed data-driven control loop.

\subsection{Privacy, Security, and Governance}
As data permeates all aspects of life, concerns of privacy and security have become paramount. Unlike oil, data can be copied, and its misuse can harm individuals or societies. Regulations like GDPR illustrate the legal response to the data economy’s challenges. Enterprises must now embed privacy by design into data platforms, and secure data against breaches. At the same time, balancing openness with control is tricky: too much restriction can stifle innovation, too little can erode trust. Data governance frameworks (including data provenance, lineage, and ethical standards) are an active area of research and policy development.

\subsection{Infrastructure and Sustainability}
The infrastructure supporting the data economy is immense. Beyond data centers, companies invest in cloud platforms, edge computing nodes, and specialized hardware (TPUs, ASICs). This infrastructure underpins AI development (the training and serving of models), as well as ubiquitous connectivity. However, the environmental footprint is rising. Researchers are exploring “green AI” and more efficient architectures to mitigate this impact. There is also a socio-economic dimension: data centers and AI growth raise questions about digital divides, national competitiveness, and the workforce needed for a data-driven future.

\subsection{Challenges and Outlook}
The data-centric world brings challenges that span technical and societal domains. From a technical standpoint, integrating disparate data sources at scale remains difficult. Model interpretability, bias, and robustness are active research areas, especially as foundation models exhibit unpredictable behavior. From a societal view, issues of algorithmic fairness, digital rights, and economic concentration (data monopolies) are pressing. Nonetheless, the potential benefits are vast: predictive medicine, optimized energy grids, intelligent transportation, and personalized education. 

In sum, data now pervades every sector. It is the raw material for AI, the basis of smart systems like digital twins, and the core of modern platforms. This pervasiveness naturally motivates the research agenda of this thesis: developing architectures and methods (from data management to AI algorithms) that harness data effectively. In the following chapters, we will delve into specific data platform architectures and digital twin methodologies that address the challenges outlined here.
